\begin{ex}
Góc nhọn có số đo bé hơn bao nhiêu độ?
\shortans{90}
\loigiai{
Bé hơn $90^\circ$.
}
\end{ex}

\begin{ex}
Góc tù có số đo lớn hơn bao nhiêu độ và bé hơn $180^\circ$?
\shortans{90}
\loigiai{
Lớn hơn $90^\circ$.
}
\end{ex}

\begin{ex}
Góc bẹt có số đo bằng bao nhiêu độ?
\shortans{180}
\loigiai{
$180^\circ$.
}
\end{ex}

\begin{ex}
Góc vuông có số đo bằng bao nhiêu độ?
\shortans{90}
\loigiai{
$90^\circ$.
}
\end{ex}

\begin{ex}
Hai đường thẳng vuông góc tạo thành mấy góc vuông?
\shortans{4}
\loigiai{
Có 4 góc vuông.
}
\end{ex}

\begin{ex}
Hai đường thẳng song song có bao nhiêu điểm chung?
\shortans{0}
\loigiai{
0 điểm chung.
}
\end{ex}

\begin{ex}
Hình chữ nhật có bao nhiêu cặp cạnh song song?
\shortans{2}
\loigiai{
Có 2 cặp.
}
\end{ex}

\begin{ex}
Hình vuông có bao nhiêu cạnh bằng nhau?
\shortans{4}
\loigiai{
Có 4 cạnh bằng nhau.
}
\end{ex}

\begin{ex}
Hình thoi có bao nhiêu cạnh bằng nhau?
\shortans{4}
\loigiai{
Có 4 cạnh bằng nhau.
}
\end{ex}

\begin{ex}
Hai đường chéo của hình thoi có vuông góc với nhau không? (Có ghi 1, không ghi 0).
\shortans{1}
\loigiai{
Có vuông góc (1).
}
\end{ex}

\begin{ex}
Một hình vuông có cạnh dài 9 cm. Chu vi hình vuông đó là bao nhiêu xăng-ti-mét?
\shortans{36}
\loigiai{
$9 \times 4 = 36$ cm.
}
\end{ex}

\begin{ex}
Một hình vuông có cạnh dài 9 cm. Diện tích hình vuông đó là bao nhiêu xăng-ti-mét vuông?
\shortans{81}
\loigiai{
$9 \times 9 = 81\text{ cm}^2$.
}
\end{ex}

\begin{ex}
Một hình chữ nhật có chiều dài 14 cm, rộng 6 cm. Chu vi là bao nhiêu xăng-ti-mét?
\shortans{40}
\loigiai{
$(14 + 6) \times 2 = 40$ cm.
}
\end{ex}

\begin{ex}
Một hình chữ nhật có chiều dài 14 cm, rộng 6 cm. Diện tích là bao nhiêu xăng-ti-mét vuông?
\shortans{84}
\loigiai{
$14 \times 6 = 84\text{ cm}^2$.
}
\end{ex}

\begin{ex}
Một hình bình hành có độ dài đáy 18 cm, chiều cao 5 cm. Diện tích là bao nhiêu xăng-ti-mét vuông?
\shortans{90}
\loigiai{
$18 \times 5 = 90\text{ cm}^2$.
}
\end{ex}

\begin{ex}
Một hình thoi có hai đường chéo dài 16 cm và 10 cm. Diện tích là bao nhiêu xăng-ti-mét vuông?
\shortans{80}
\loigiai{
$16 \times 10 : 2 = 80\text{ cm}^2$.
}
\end{ex}

\begin{ex}
Hình thoi có chu vi 52 cm. Cạnh của hình thoi dài bao nhiêu xăng-ti-mét?
\shortans{13}
\loigiai{
$52 : 4 = 13$ cm.
}
\end{ex}

\begin{ex}
Lúc 6 giờ đúng, góc giữa kim giờ và kim phút bằng bao nhiêu độ?
\shortans{180}
\loigiai{
$180^\circ$.
}
\end{ex}

\begin{ex}
Lúc 3 giờ đúng, góc giữa hai kim bằng bao nhiêu độ?
\shortans{90}
\loigiai{
$90^\circ$.
}
\end{ex}

\begin{ex}
Lúc 2 giờ đúng, góc giữa hai kim bằng bao nhiêu độ?
\shortans{60}
\loigiai{
$30 \times 2 = 60^\circ$.
}
\end{ex}

\begin{ex}
Hình thang có mấy cặp cạnh đáy song song?
\shortans{1}
\loigiai{
Có 1 cặp.
}
\end{ex}

\begin{ex}
Một hình bình hành có chu vi 48 cm, một cạnh dài 15 cm. Cạnh kề dài bao nhiêu xăng-ti-mét?
\shortans{9}
\loigiai{
Nửa chu vi = 24 cm. Cạnh kia = $24 - 15 = 9$ cm.
}
\end{ex}

\begin{ex}
Một hình vuông có diện tích $64\text{ cm}^2$. Chu vi hình vuông đó là bao nhiêu xăng-ti-mét?
\shortans{32}
\loigiai{
Cạnh = 8 cm. Chu vi = $8 \times 4 = 32$ cm.
}
\end{ex}

\begin{ex}
Hình chữ nhật có chu vi 44 cm, chiều dài 15 cm. Diện tích hình chữ nhật đó là bao nhiêu xăng-ti-mét vuông?
\shortans{105}
\loigiai{
Rộng = 7 cm. Diện tích = $15 \times 7 = 105\text{ cm}^2$.
}
\end{ex}

\begin{ex}
Một hình tam giác có 1 góc vuông thì gọi là tam giác gì? (Tam giác vuông ghi 1, tam giác tù ghi 2).
\shortans{1}
\loigiai{
Tam giác vuông (1).
}
\end{ex}
